%!TEX root = ../../thesis.tex

\subsection{cJSON}\label{sec:target-cjson}

% TODO: The first target whose reward is not designed around the agent, and the only
% one run under both formulations --- state both facts in the opening paragraph.
%
% The target: the cJSON parser, driven through the emulator by a harness that reads a
% fixed-size record and parses it. Describe the two harness styles actually used and
% why they differ:
%   - the input-level environment re-parses the accumulated prefix after every byte
%     (cJSON has no streaming API), so each step yields both a coverage observation of
%     the current parse state and a validity verdict,
%   - the corpus environment uses a persistent-mode loop that reads one whole record
%     per execution and never exits, which is what keeps the boot snapshot restorable
%     across a sweep.
% Note that the harness is part of the experimental apparatus: a target is a program
% *plus* a harness, and the harness differs between the two formulations by necessity.
% Any coverage comparison between them must account for that
% (\cref{sec:results-head-to-head}).
%
% What it tests: whether coverage alone carries enough signal to drive a parser into
% its grammar --- valid JSON reaches deep parser branches, malformed input bails out
% early. Report the measured coverage separation between valid and malformed input as
% evidence that the signal exists at all, independent of any learning result.
%
% Action alphabet: the restricted JSON-relevant set for the input-level formulation
% (\cref{sec:input-action-space}) versus byte-level operators for the corpus
% formulation (\cref{sec:corpus-action-space}). This asymmetry is the main confound in
% the head-to-head and must be named here.
